\section{Use Case Descriptions}

\subsection*{Module 1: User Management System}
\label{module:minipage:user_management}

% --- UC-1: User Registration ---
\begin{minipage}{\textwidth}
\begin{center}
    \textbf{Use Case: User Registration (UC-1)}
\end{center}
\vspace{0.5ex}
\small

\noindent
\begin{tabularx}{\textwidth}{|p{0.28\textwidth}|X|}
    \hline
    \textbf{Property} & \textbf{Details} \\
    \hline
    Use Case Name & User Registration \\
    \hline
    Use Case ID & UC-1 \\
    \hline
    Description & This use case allows new users to create an account on the Verdant Search platform by providing registration details. \\
    \hline
    Priority & High \\
    \hline
    Actor & User (prospective), System \\
    \hline
    Pre-condition &
        \parbox[t]{\linewidth}{%
        \setlist[itemize]{nosep, leftmargin=1em, itemsep=0pt, parsep=0pt, topsep=0pt, partopsep=0pt, after=\strut}
        \begin{itemize}
            \item User does not have an existing account in the system.
            \item System is operational and accessible.
        \end{itemize}%
        } \\
    \hline
    Post-condition &
        \parbox[t]{\linewidth}{%
        \setlist[itemize]{nosep, leftmargin=1em, itemsep=0pt, parsep=0pt, topsep=0pt, partopsep=0pt, after=\strut}
        \begin{itemize}
            \item A new user account is created and stored in the database.
            \item Verification email is sent to the registered email address.
            \item User can proceed to login after email verification.
        \end{itemize}%
        } \\
    \hline
    Triggering event &
        \parbox[t]{\linewidth}{%
        A prospective user navigates to the registration page and submits registration details.
        } \\
    \hline
    Flow of Events (Main Success Scenario) &
        \parbox[t]{\linewidth}{%
        \setlist[enumerate]{nosep, leftmargin=1.5em, itemsep=0pt, parsep=0pt, topsep=0pt, partopsep=0pt, after=\strut}
        \begin{enumerate}
            \item User navigates to the registration page.
            \item System displays registration form requesting email and password.
            \item User enters email address and chooses a password (FR-1.1).
            \item System validates email format and password strength requirements (minimum 8 characters, uppercase, lowercase, number, special character) (FR-1.2).
            \item System checks if email already exists in database.
            \item System creates new user account with hashed password.
            \item System sends verification email to registered address (FR-1.3).
            \item System displays confirmation message instructing user to verify email.
        \end{enumerate}%
        } \\
    \hline
    Alternative Flow &
        \parbox[t]{\linewidth}{%
        \setlist[itemize]{nosep, leftmargin=1em, itemsep=0pt, parsep=0pt, topsep=0pt, partopsep=0pt, after=\strut}
        \begin{itemize}
            \item User clicks verification link in email: System marks account as verified and redirects to login page.
        \end{itemize}%
        } \\
    \hline
    Exception Flow &
        \parbox[t]{\linewidth}{%
        \setlist[itemize]{nosep, leftmargin=1em, itemsep=0pt, parsep=0pt, topsep=0pt, partopsep=0pt, after=\strut}
        \begin{itemize}
            \item If email already exists: System displays error message "Email already registered" and prompts user to login or reset password.
            \item If password does not meet strength requirements: System displays specific validation error (e.g., "Password must contain at least one uppercase letter").
            \item If email format is invalid: System displays error "Please enter a valid email address".
            \item If email service fails: System logs error, creates account but displays warning that verification email may be delayed.
        \end{itemize}%
        } \\
    \hline
\end{tabularx}
\end{minipage}
\bigskip

% --- UC-2: User Login via SSO ---
\begin{minipage}{\textwidth}
\begin{center}
    \textbf{Use Case: User Login via SSO (UC-2)}
\end{center}
\vspace{0.5ex}
\small

\noindent
\begin{tabularx}{\textwidth}{|p{0.28\textwidth}|X|}
    \hline
    \textbf{Property} & \textbf{Details} \\
    \hline
    Use Case Name & User Login via SSO \\
    \hline
    Use Case ID & UC-2 \\
    \hline
    Description & This use case allows users to authenticate using Single Sign-On providers (LDAP, Okta, Azure AD) integrated with Infinity Datatech's identity management system. \\
    \hline
    Priority & High \\
    \hline
    Actor & User, System, SSO Provider \\
    \hline
    Pre-condition &
        \parbox[t]{\linewidth}{%
        \setlist[itemize]{nosep, leftmargin=1em, itemsep=0pt, parsep=0pt, topsep=0pt, partopsep=0pt, after=\strut}
        \begin{itemize}
            \item User has valid credentials with the SSO provider.
            \item System has configured SSO integration with the provider.
            \item System is operational and reachable.
        \end{itemize}%
        } \\
    \hline
    Post-condition &
        \parbox[t]{\linewidth}{%
        \setlist[itemize]{nosep, leftmargin=1em, itemsep=0pt, parsep=0pt, topsep=0pt, partopsep=0pt, after=\strut}
        \begin{itemize}
            \item User is authenticated and session is established.
            \item JWT token is generated and provided to user client.
            \item User is redirected to main search interface.
        \end{itemize}%
        } \\
    \hline
    Triggering event &
        \parbox[t]{\linewidth}{%
        User clicks "Sign in with SSO" button on login page.
        } \\
    \hline
    Flow of Events (Main Success Scenario) &
        \parbox[t]{\linewidth}{%
        \setlist[enumerate]{nosep, leftmargin=1.5em, itemsep=0pt, parsep=0pt, topsep=0pt, partopsep=0pt, after=\strut}
        \begin{enumerate}
            \item User navigates to login page and selects SSO authentication option.
            \item System redirects user to configured SSO provider login page.
            \item User authenticates with SSO provider using their corporate credentials.
            \item SSO provider validates credentials using LDAP, Okta, or Azure AD (FR-1.4).
            \item SSO provider returns authentication token using OAuth2 or SAML protocol (FR-1.5).
            \item System validates authentication token from SSO provider.
            \item System creates or retrieves user account based on SSO identifier.
            \item System generates JWT token for session management (FR-1.6).
            \item System stores JWT token in secure HTTP-only cookie.
            \item System redirects user to main search interface.
        \end{enumerate}%
        } \\
    \hline
    Alternative Flow &
        \parbox[t]{\linewidth}{%
        \setlist[itemize]{nosep, leftmargin=1em, itemsep=0pt, parsep=0pt, topsep=0pt, partopsep=0pt, after=\strut}
        \begin{itemize}
            \item First-time SSO user: System creates new account automatically using SSO profile information.
        \end{itemize}%
        } \\
    \hline
    Exception Flow &
        \parbox[t]{\linewidth}{%
        \setlist[itemize]{nosep, leftmargin=1em, itemsep=0pt, parsep=0pt, topsep=0pt, partopsep=0pt, after=\strut}
        \begin{itemize}
            \item If SSO authentication fails: System displays error message and redirects back to login page.
            \item If SSO provider is unreachable: System displays service unavailable message and suggests using email/password login.
            \item If authentication token validation fails: System logs security event and denies access.
            \item If user account is disabled: System displays "Account suspended" message and contacts administrator.
        \end{itemize}%
        } \\
    \hline
\end{tabularx}
\end{minipage}
\bigskip

% --- UC-12: Submit Search Query ---
\begin{minipage}{\textwidth}
\begin{center}
    \textbf{Use Case: Submit Search Query (UC-12)}
\end{center}
\vspace{0.5ex}
\small

\noindent
\begin{tabularx}{\textwidth}{|p{0.28\textwidth}|X|}
    \hline
    \textbf{Property} & \textbf{Details} \\
    \hline
    Use Case Name & Submit Search Query \\
    \hline
    Use Case ID & UC-12 \\
    \hline
    Description & This use case allows users to initiate a search by submitting a text query, which triggers the hybrid retrieval pipeline combining BM25 and vector search. \\
    \hline
    Priority & High \\
    \hline
    Actor & User, System \\
    \hline
    Pre-condition &
        \parbox[t]{\linewidth}{%
        \setlist[itemize]{nosep, leftmargin=1em, itemsep=0pt, parsep=0pt, topsep=0pt, partopsep=0pt, after=\strut}
        \begin{itemize}
            \item User is authenticated and has active session.
            \item Search indices (BM25 inverted index and HNSW vector index) are built and available.
            \item System backend services are operational.
        \end{itemize}%
        } \\
    \hline
    Post-condition &
        \parbox[t]{\linewidth}{%
        \setlist[itemize]{nosep, leftmargin=1em, itemsep=0pt, parsep=0pt, topsep=0pt, partopsep=0pt, after=\strut}
        \begin{itemize}
            \item Query is processed and sanitized.
            \item Search query is logged for analytics.
            \item Hybrid retrieval pipeline is triggered.
            \item User receives ranked search results.
        \end{itemize}%
        } \\
    \hline
    Triggering event &
        \parbox[t]{\linewidth}{%
        User types query into search box and presses Enter or clicks Search button.
        } \\
    \hline
    Flow of Events (Main Success Scenario) &
        \parbox[t]{\linewidth}{%
        \setlist[enumerate]{nosep, leftmargin=1.5em, itemsep=0pt, parsep=0pt, topsep=0pt, partopsep=0pt, after=\strut}
        \begin{enumerate}
            \item User enters search query in text input field (FR-2.1).
            \item System validates query is not empty and within length limits.
            \item System sanitizes query to prevent injection attacks (removes special characters, validates encoding) (FR-2.2).
            \item System normalizes query text (lowercasing, removing extra whitespace).
            \item System optionally expands query with synonyms if synonym detection is enabled (FR-2.3).
            \item System logs query with user ID, timestamp, and session information.
            \item System initiates parallel execution of BM25 keyword search (UC-16) and vector search (UC-17).
            \item System displays loading indicator to user.
        \end{enumerate}%
        } \\
    \hline
    Alternative Flow &
        \parbox[t]{\linewidth}{%
        \setlist[itemize]{nosep, leftmargin=1em, itemsep=0pt, parsep=0pt, topsep=0pt, partopsep=0pt, after=\strut}
        \begin{itemize}
            \item User refines existing search: System preserves previous context and stores query history.
            \item System detects similar recent query: System retrieves cached results if available within freshness window (5 minutes).
        \end{itemize}%
        } \\
    \hline
    Exception Flow &
        \parbox[t]{\linewidth}{%
        \setlist[itemize]{nosep, leftmargin=1em, itemsep=0pt, parsep=0pt, topsep=0pt, partopsep=0pt, after=\strut}
        \begin{itemize}
            \item If query is empty: System displays "Please enter a search query" message.
            \item If query exceeds maximum length (1000 characters): System displays "Query too long, please shorten" error.
            \item If sanitization detects potential injection attack: System logs security event and displays generic error.
            \item If backend search service is unavailable: System displays "Search service temporarily unavailable" and suggests retry.
        \end{itemize}%
        } \\
    \hline
\end{tabularx}
\end{minipage}
\bigskip


% --- UC-16: Perform Keyword Search ---
\begin{minipage}{\textwidth}
\begin{center}
    \textbf{Use Case: Perform Keyword Search (UC-16)}
\end{center}
\vspace{0.5ex}
\small

\noindent
\begin{tabularx}{\textwidth}{|p{0.28\textwidth}|X|}
    \hline
    \textbf{Property} & \textbf{Details} \\
    \hline
    Use Case Name & Perform Keyword Search \\
    \hline
    Use Case ID & UC-16 \\
    \hline
    Description & This use case executes BM25 keyword-based retrieval using the inverted index to identify documents containing query terms. \\
    \hline
    Priority & High \\
    \hline
    Actor & System \\
    \hline
    Pre-condition &
        \parbox[t]{\linewidth}{%
        \setlist[itemize]{nosep, leftmargin=1em, itemsep=0pt, parsep=0pt, topsep=0pt, partopsep=0pt, after=\strut}
        \begin{itemize}
            \item Sanitized and normalized query is available from UC-12.
            \item BM25 inverted index has been built and is loaded in memory.
            \item Index contains document term frequencies and collection statistics.
        \end{itemize}%
        } \\
    \hline
    Post-condition &
        \parbox[t]{\linewidth}{%
        \setlist[itemize]{nosep, leftmargin=1em, itemsep=0pt, parsep=0pt, topsep=0pt, partopsep=0pt, after=\strut}
        \begin{itemize}
            \item Top K documents are retrieved based on BM25 scores.
            \item Each result includes document ID, score, and matching term positions.
            \item Results are passed to hybrid fusion component (UC-18).
        \end{itemize}%
        } \\
    \hline
    Triggering event &
        \parbox[t]{\linewidth}{%
        Query submission (UC-12) triggers parallel keyword and vector search.
        } \\
    \hline
    Flow of Events (Main Success Scenario) &
        \parbox[t]{\linewidth}{%
        \setlist[enumerate]{nosep, leftmargin=1.5em, itemsep=0pt, parsep=0pt, topsep=0pt, partopsep=0pt, after=\strut}
        \begin{enumerate}
            \item System tokenizes query into individual terms.
            \item System looks up each term in inverted index to retrieve posting lists (FR-2.13).
            \item For each document in posting lists, system computes BM25 score using formula:
            $$BM25(D,Q) = \sum_{t \in Q} IDF(t) \cdot \frac{f(t,D) \cdot (k_1 + 1)}{f(t,D) + k_1 \cdot (1 - b + b \cdot \frac{|D|}{avgdl})}$$
            where $k_1=1.2$ and $b=0.75$ are tunable parameters (FR-2.14).
            \item System ranks documents by descending BM25 score.
            \item System retrieves top K=100 candidates with their scores and term positions (FR-2.15).
            \item System returns results to hybrid search fusion component.
        \end{enumerate}%
        } \\
    \hline
    Alternative Flow &
        \parbox[t]{\linewidth}{%
        \setlist[itemize]{nosep, leftmargin=1em, itemsep=0pt, parsep=0pt, topsep=0pt, partopsep=0pt, after=\strut}
        \begin{itemize}
            \item If query contains stop words: System optionally removes them based on configuration.
            \item If query contains phrase: System uses term position information for proximity scoring.
        \end{itemize}%
        } \\
    \hline
    Exception Flow &
        \parbox[t]{\linewidth}{%
        \setlist[itemize]{nosep, leftmargin=1em, itemsep=0pt, parsep=0pt, topsep=0pt, partopsep=0pt, after=\strut}
        \begin{itemize}
            \item If no documents match any query terms: System returns empty result set with message "No documents found".
            \item If inverted index is corrupted or unavailable: System logs critical error and falls back to vector-only search.
            \item If BM25 computation exceeds time limit (1 second): System returns partial results and logs performance warning.
        \end{itemize}%
        } \\
    \hline
\end{tabularx}
\end{minipage}
\bigskip

% --- UC-20: Apply Multimodal Reranking ---
\begin{minipage}{\textwidth}
\begin{center}
    \textbf{Use Case: Apply Multimodal Reranking (UC-20)}
\end{center}
\vspace{0.5ex}
\small

\noindent
\begin{tabularx}{\textwidth}{|p{0.28\textwidth}|X|}
    \hline
    \textbf{Property} & \textbf{Details} \\
    \hline
    Use Case Name & Apply Multimodal Reranking \\
    \hline
    Use Case ID & UC-20 \\
    \hline
    Description & This use case applies a BLIP-2 based multimodal listwise reranker to top candidates from hybrid search, evaluating both textual content and visual layouts for relevance. \\
    \hline
    Priority & High \\
    \hline
    Actor & System \\
    \hline
    Pre-condition &
        \parbox[t]{\linewidth}{%
        \setlist[itemize]{nosep, leftmargin=1em, itemsep=0pt, parsep=0pt, topsep=0pt, partopsep=0pt, after=\strut}
        \begin{itemize}
            \item Hybrid search (UC-18) has produced merged top K=20 candidates.
            \item BLIP-2 multimodal reranking model is loaded in GPU memory.
            \item Document texts and associated images are accessible.
        \end{itemize}%
        } \\
    \hline
    Post-condition &
        \parbox[t]{\linewidth}{%
        \setlist[itemize]{nosep, leftmargin=1em, itemsep=0pt, parsep=0pt, topsep=0pt, partopsep=0pt, after=\strut}
        \begin{itemize}
            \item Documents are reranked based on multimodal relevance scores.
            \item Final ranked list is returned to user interface.
            \item Reranking latency is logged for performance monitoring.
        \end{itemize}%
        } \\
    \hline
    Triggering event &
        \parbox[t]{\linewidth}{%
        Hybrid search fusion completes and provides top K candidates.
        } \\
    \hline
    Flow of Events (Main Success Scenario) &
        \parbox[t]{\linewidth}{%
        \setlist[enumerate]{nosep, leftmargin=1.5em, itemsep=0pt, parsep=0pt, topsep=0pt, partopsep=0pt, after=\strut}
        \begin{enumerate}
            \item System receives top K=20 candidate documents from hybrid search (FR-2.25).
            \item For each candidate, system retrieves document text and associated images.
            \item System constructs multimodal prompt following setwise prompting format:
            \begin{verbatim}
Query: [user query]
Documents:
[1] Text: ... Images: [image1, image2]
[2] Text: ... Images: [image3]
...
Rank these documents by relevance.
            \end{verbatim}
            \item System encodes prompt and images using BLIP-2 vision-language model (FR-2.26).
            \item BLIP-2 Q-Former processes cross-modal interactions between text and images.
            \item System decodes listwise ranking from model output.
            \item System reorders candidates based on multimodal relevance scores (FR-2.27).
            \item System returns final ranked list to display component.
        \end{enumerate}%
        } \\
    \hline
    Alternative Flow &
        \parbox[t]{\linewidth}{%
        \setlist[itemize]{nosep, leftmargin=1em, itemsep=0pt, parsep=0pt, topsep=0pt, partopsep=0pt, after=\strut}
        \begin{itemize}
            \item Document has no images: System applies text-only reranking using BLIP-2 text encoder.
            \item Candidate list has fewer than 20 documents: System reranks available candidates.
        \end{itemize}%
        } \\
    \hline
    Exception Flow &
        \parbox[t]{\linewidth}{%
        \setlist[itemize]{nosep, leftmargin=1em, itemsep=0pt, parsep=0pt, topsep=0pt, partopsep=0pt, after=\strut}
        \begin{itemize}
            \item If BLIP-2 model inference fails: System logs error and returns original hybrid search ranking.
            \item If reranking exceeds time limit (3 seconds): System terminates reranking and uses partial results.
            \item If GPU memory is exhausted: System falls back to CPU inference or skips reranking.
            \item If image loading fails: System excludes failed images and continues with available visual content.
        \end{itemize}%
        } \\
    \hline
\end{tabularx}
\end{minipage}
\bigskip

% --- UC-26: Trigger Distributed Crawling ---
\begin{minipage}{\textwidth}
\begin{center}
    \textbf{Use Case: Trigger Distributed Crawling (UC-26)}
\end{center}
\vspace{0.5ex}
\small

\noindent
\begin{tabularx}{\textwidth}{|p{0.28\textwidth}|X|}
    \hline
    \textbf{Property} & \textbf{Details} \\
    \hline
    Use Case Name & Trigger Distributed Crawling \\
    \hline
    Use Case ID & UC-26 \\
    \hline
    Description & This use case allows administrators to initiate distributed crawling jobs coordinated via Kafka to ingest documents from configured data sources. \\
    \hline
    Priority & High \\
    \hline
    Actor & Administrator, System \\
    \hline
    Pre-condition &
        \parbox[t]{\linewidth}{%
        \setlist[itemize]{nosep, leftmargin=1em, itemsep=0pt, parsep=0pt, topsep=0pt, partopsep=0pt, after=\strut}
        \begin{itemize}
            \item Administrator is authenticated with appropriate permissions.
            \item Data sources are configured with valid credentials (UC-25).
            \item Kafka cluster and crawler worker nodes are operational.
        \end{itemize}%
        } \\
    \hline
    Post-condition &
        \parbox[t]{\linewidth}{%
        \setlist[itemize]{nosep, leftmargin=1em, itemsep=0pt, parsep=0pt, topsep=0pt, partopsep=0pt, after=\strut}
        \begin{itemize}
            \item Crawling job is created and queued in Kafka.
            \item Worker nodes begin processing crawl tasks.
            \item Crawling progress is visible in monitoring dashboard.
            \item Documents are extracted and queued for indexing.
        \end{itemize}%
        } \\
    \hline
    Triggering event &
        \parbox[t]{\linewidth}{%
        Administrator clicks "Start Crawl" button in admin dashboard or scheduled crawl time arrives.
        } \\
    \hline
    Flow of Events (Main Success Scenario) &
        \parbox[t]{\linewidth}{%
        \setlist[enumerate]{nosep, leftmargin=1.5em, itemsep=0pt, parsep=0pt, topsep=0pt, partopsep=0pt, after=\strut}
        \begin{enumerate}
            \item Administrator selects data sources to crawl from configuration list.
            \item Administrator clicks "Start Crawl" button (FR-3.4).
            \item System creates crawl job with unique job ID and timestamp.
            \item System partitions crawl tasks across configured data sources.
            \item System publishes crawl tasks to Kafka topic "crawl-tasks" (FR-3.5).
            \item Kafka distributes tasks to available crawler worker nodes.
            \item Each worker node consumes assigned tasks and begins crawling.
            \item Workers report progress back to coordinator via Kafka "crawl-status" topic.
            \item System aggregates progress metrics and updates dashboard.
            \item Upon completion, workers publish extracted documents to "documents-ingested" topic.
        \end{enumerate}%
        } \\
    \hline
    Alternative Flow &
        \parbox[t]{\linewidth}{%
        \setlist[itemize]{nosep, leftmargin=1em, itemsep=0pt, parsep=0pt, topsep=0pt, partopsep=0pt, after=\strut}
        \begin{itemize}
            \item Incremental crawl: System retrieves last crawl timestamp and configures workers to fetch only updated documents.
            \item Scheduled crawl: System triggers crawl automatically at configured interval without administrator action.
        \end{itemize}%
        } \\
    \hline
    Exception Flow &
        \parbox[t]{\linewidth}{%
        \setlist[itemize]{nosep, leftmargin=1em, itemsep=0pt, parsep=0pt, topsep=0pt, partopsep=0pt, after=\strut}
        \begin{itemize}
            \item If worker node fails during crawl: System detects failure via missed heartbeats, reassigns tasks to healthy nodes (FR-3.6).
            \item If data source is unreachable: Worker implements exponential backoff (1s, 2s, 4s, 8s) for retries, then marks task as failed.
            \item If Kafka cluster is unavailable: System displays error message and prevents crawl initiation until Kafka is restored.
            \item If all workers are busy: Tasks remain queued in Kafka until workers become available.
        \end{itemize}%
        } \\
    \hline
\end{tabularx}
\end{minipage}
\bigskip

% --- UC-43: Receive User Query (RAG) ---
\begin{minipage}{\textwidth}
\begin{center}
    \textbf{Use Case: Receive User Query for RAG (UC-43)}
\end{center}
\vspace{0.5ex}
\small

\noindent
\begin{tabularx}{\textwidth}{|p{0.28\textwidth}|X|}
    \hline
    \textbf{Property} & \textbf{Details} \\
    \hline
    Use Case Name & Receive User Query for RAG \\
    \hline
    Use Case ID & UC-43 \\
    \hline
    Description & This use case handles user queries submitted to the RAG (Retrieval-Augmented Generation) system for AI-powered answer generation based on retrieved context. \\
    \hline
    Priority & High \\
    \hline
    Actor & User, System \\
    \hline
    Pre-condition &
        \parbox[t]{\linewidth}{%
        \setlist[itemize]{nosep, leftmargin=1em, itemsep=0pt, parsep=0pt, topsep=0pt, partopsep=0pt, after=\strut}
        \begin{itemize}
            \item User is authenticated with active session.
            \item Search and retrieval system is operational.
            \item LLM inference service (SGLang) is running and responsive.
        \end{itemize}%
        } \\
    \hline
    Post-condition &
        \parbox[t]{\linewidth}{%
        \setlist[itemize]{nosep, leftmargin=1em, itemsep=0pt, parsep=0pt, topsep=0pt, partopsep=0pt, after=\strut}
        \begin{itemize}
            \item Query is validated and accepted.
            \item Query is added to user conversation history.
            \item Retrieval pipeline is triggered to fetch relevant context.
            \item Query is queued for LLM processing.
        \end{itemize}%
        } \\
    \hline
    Triggering event &
        \parbox[t]{\linewidth}{%
        User submits question in RAG chat interface and clicks Send or presses Enter.
        } \\
    \hline
    Flow of Events (Main Success Scenario) &
        \parbox[t]{\linewidth}{%
        \setlist[enumerate]{nosep, leftmargin=1.5em, itemsep=0pt, parsep=0pt, topsep=0pt, partopsep=0pt, after=\strut}
        \begin{enumerate}
            \item User types question into RAG chat input field.
            \item User submits query by clicking Send button or pressing Enter.
            \item System receives query via WebSocket or HTTP POST request (FR-5.1).
            \item System validates query length is within 1000 character limit (FR-5.2).
            \item System checks current system load and available LLM capacity.
            \item System adds query to user conversation history stored in Redis.
            \item System triggers hybrid retrieval pipeline (UC-18) to fetch relevant document chunks.
            \item System displays typing indicator to user while processing.
            \item Upon retrieval completion, system proceeds to context injection (UC-47).
        \end{enumerate}%
        } \\
    \hline
    Alternative Flow &
        \parbox[t]{\linewidth}{%
        \setlist[itemize]{nosep, leftmargin=1em, itemsep=0pt, parsep=0pt, topsep=0pt, partopsep=0pt, after=\strut}
        \begin{itemize}
            \item Follow-up question detected: System identifies query as follow-up (UC-45) and triggers query rewriting (UC-55) before retrieval.
        \end{itemize}%
        } \\
    \hline
    Exception Flow &
        \parbox[t]{\linewidth}{%
        \setlist[itemize]{nosep, leftmargin=1em, itemsep=0pt, parsep=0pt, topsep=0pt, partopsep=0pt, after=\strut}
        \begin{itemize}
            \item If query exceeds 1000 characters: System displays "Query too long, please shorten to 1000 characters" error.
            \item If system is under high load: System queues query (FR-5.3) and displays "Your question is queued, estimated wait: X seconds".
            \item If LLM service is unavailable: System displays "Answer generation temporarily unavailable, please try again" message.
            \item If retrieval returns no results: System proceeds to answer generation (UC-54) which indicates knowledge gap.
        \end{itemize}%
        } \\
    \hline
\end{tabularx}
\end{minipage}
\bigskip

% --- UC-51: Generate Grounded Answer ---
\begin{minipage}{\textwidth}
\begin{center}
    \textbf{Use Case: Generate Grounded Answer (UC-51)}
\end{center}
\vspace{0.5ex}
\small

\noindent
\begin{tabularx}{\textwidth}{|p{0.28\textwidth}|X|}
    \hline
    \textbf{Property} & \textbf{Details} \\
    \hline
    Use Case Name & Generate Grounded Answer \\
    \hline
    Use Case ID & UC-51 \\
    \hline
    Description & This use case controls LLM answer generation to ensure responses are grounded exclusively in retrieved context, preventing hallucination. \\
    \hline
    Priority & High \\
    \hline
    Actor & System, LLM \\
    \hline
    Pre-condition &
        \parbox[t]{\linewidth}{%
        \setlist[itemize]{nosep, leftmargin=1em, itemsep=0pt, parsep=0pt, topsep=0pt, partopsep=0pt, after=\strut}
        \begin{itemize}
            \item Context has been injected into LLM prompt (UC-47).
            \item LLM model is loaded and ready for inference.
            \item System prompt instructs model to avoid hallucination.
        \end{itemize}%
        } \\
    \hline
    Post-condition &
        \parbox[t]{\linewidth}{%
        \setlist[itemize]{nosep, leftmargin=1em, itemsep=0pt, parsep=0pt, topsep=0pt, partopsep=0pt, after=\strut}
        \begin{itemize}
            \item LLM generates answer based only on provided context.
            \item Answer includes inline citations to source chunks.
            \item Uncertain claims are marked with confidence qualifiers.
            \item Generated answer is streamed to user (UC-48).
        \end{itemize}%
        } \\
    \hline
    Triggering event &
        \parbox[t]{\linewidth}{%
        Context injection completes and final prompt is ready for LLM inference.
        } \\
    \hline
    Flow of Events (Main Success Scenario) &
        \parbox[t]{\linewidth}{%
        \setlist[enumerate]{nosep, leftmargin=1.5em, itemsep=0pt, parsep=0pt, topsep=0pt, partopsep=0pt, after=\strut}
        \begin{enumerate}
            \item System constructs system prompt instructing LLM:
            \begin{verbatim}
You are an expert search assistant. Answer the question
using ONLY the provided context. If the context doesn't
contain enough information, say so clearly. Always cite
sources using [1], [2] notation.
            \end{verbatim}
            (FR-5.26)
            \item System appends retrieved context chunks with source identifiers.
            \item System appends user question to prompt.
            \item System submits complete prompt to LLM via SGLang inference API.
            \item LLM generates answer tokens using only provided context (FR-5.25).
            \item LLM includes citation markers [1], [2] referencing specific context chunks.
            \item For uncertain claims, LLM uses qualifiers like "According to the documentation..." or "The context suggests..." (FR-5.27).
            \item System streams generated tokens to user interface (UC-48).
        \end{enumerate}%
        } \\
    \hline
    Alternative Flow &
        \parbox[t]{\linewidth}{%
        \setlist[itemize]{nosep, leftmargin=1em, itemsep=0pt, parsep=0pt, topsep=0pt, partopsep=0pt, after=\strut}
        \begin{itemize}
            \item Context insufficient: LLM generates message "I cannot find enough information in the knowledge base to answer this question" (handled by UC-54).
        \end{itemize}%
        } \\
    \hline
    Exception Flow &
        \parbox[t]{\linewidth}{%
        \setlist[itemize]{nosep, leftmargin=1em, itemsep=0pt, parsep=0pt, topsep=0pt, partopsep=0pt, after=\strut}
        \begin{itemize}
            \item If LLM generates unsupported claim: Entailment validation (UC-53) flags the claim for review.
            \item If LLM fails to include citations: System post-processes answer to add citations based on context matching.
            \item If generation exceeds time limit (30 seconds): System terminates generation and displays partial answer with warning.
            \item If LLM produces empty response: System displays "Unable to generate answer, please rephrase question".
        \end{itemize}%
        } \\
    \hline
\end{tabularx}
\end{minipage}
\bigskip

\clearpage
